\documentclass[12pt]{article}
\usepackage[utf8]{inputenc}
\usepackage{amsmath, amssymb}
\usepackage{xcolor}
\usepackage{geometry}
\usepackage{hyperref}
\usepackage{fancyhdr}
\usepackage{enumitem}
\usepackage{minted}
\usepackage{booktabs}
\usepackage{tikz}
\usetikzlibrary{shapes, arrows, positioning}

\geometry{margin=1in}
\hypersetup{colorlinks=true, linkcolor=blue, urlcolor=cyan}

\pagestyle{fancy}
\fancyhf{}
\fancyhead[L]{\textbf{\TOPICTITLE}}
\fancyhead[R]{\thepage}

% ------------------------------- 
% Topic Metadata
% ------------------------------- 
\newcommand{\TOPICTITLE}{Link Layer}
\title{\TOPICTITLE\\\large Study-Ready Notes}
\author{Compiled by Andrew Photinakis}
\date{\today}

\setlength{\headheight}{15pt}

\begin{document}
\maketitle
\tableofcontents
\newpage

\section*{Chapter 6 - Link Layer}

\section{Introduction to the Link Layer (Section 6.1)}

\subsection{Conceptual Framework: The Node-to-Node Perspective}

\subsubsection{Core Terminology and Abstraction}

While the network layer (IP) is responsible for routing a datagram from the \emph{source} host to the \emph{final destination}, the link layer is concerned with a more localized task: moving that datagram across a \textbf{single physical link} between two adjacent nodes.

To describe link-layer behavior precisely, we introduce several key terms:

\begin{itemize}
    \item \textbf{Nodes}: Devices that implement link-layer protocols (hosts, routers, switches, WiFi access points).
    \item \textbf{Links}: Communication channels connecting adjacent nodes. These may be wired (copper, fiber) or wireless (WiFi, satellite).
    \item \textbf{Frame}: The link-layer packet. A network-layer datagram is encapsulated into a frame before transmission.
\end{itemize}

\paragraph{Heterogeneity of Links.}
A fundamental concept is that as a datagram travels end-to-end, it typically traverses \emph{multiple, heterogeneous link types}. Each individual link may provide different service guarantees.

\begin{itemize}
    \item A datagram might begin its journey on an Ethernet LAN, traverse a fiber backbone, and conclude on a WiFi hop.
    \item Each link-layer protocol independently determines its own reliability and error-handling procedures.
\end{itemize}

\noindent The implication is that link-layer protocols are \emph{not uniform} across the network path; each operates autonomously on its specific link.

% ----------------------------------------------------------------------
\subsection{The Transportation Analogy}

To develop intuition for the relationship between the network and link layers, consider a tourist traveling from Princeton, NJ to Lausanne, Switzerland.

\paragraph{Scenario Overview.}
A tourist (representing the \textbf{datagram}) wishes to travel internationally. A travel agent (the \textbf{routing protocol}) plans the itinerary:

\begin{enumerate}
    \item Limousine from Princeton to JFK Airport
    \item Plane from JFK to Geneva
    \item Train from Geneva to Lausanne
\end{enumerate}

\paragraph{Mapping to Networking Concepts.}

\begin{center}
    \begin{tabular}{|l|l|p{7cm}|}
        \hline
        \textbf{Transportation Analogy}          & \textbf{Networking Concept} & \textbf{Function}                       \\ \hline
        Tourist                                  & Datagram                    & Payload being transported               \\ \hline
        Travel Agent                             & Routing Protocol            & Determines end-to-end path              \\ \hline
        Leg of Trip                              & Link                        & Adjacent-node physical connection       \\ \hline
        Transportation Mode (Limo, Plane, Train) & Link-Layer Protocol         & Controls per-link transmission behavior \\ \hline
    \end{tabular}
\end{center}

\noindent \textbf{Key Insight:}
Just as the tourist is successively encapsulated within different transportation modes, the datagram is encapsulated in different link-layer frames as it moves across different links. Each protocol operates independently on its own ``leg’’ of the journey.

% ----------------------------------------------------------------------
\subsection{Services Provided by the Link Layer}

The link layer supplies a set of services to the network layer. The specific services depend on the protocol used on a particular link.

\subsubsection{Framing}

\begin{itemize}
    \item \textbf{Definition}: Encapsulation of a datagram inside a link-layer frame.
    \item \textbf{Mechanism}: A header and (often) a trailer are added:
          \[
              \text{Frame} = [\text{Header} \mid \text{Datagram} \mid \text{Trailer}]
          \]
    \item \textbf{Intuition}: Similar to enclosing a letter in an envelope with a standardized address format.
\end{itemize}

\subsubsection{Link Access (Medium Access Control, MAC)}

\begin{itemize}
    \item \textbf{Definition}: Rules governing which node transmits on a shared medium.
    \item \textbf{Point-to-Point Links}: Simple protocols where one sender communicates with one receiver.
    \item \textbf{Broadcast Links}: Require \textbf{MAC protocols} (e.g., Ethernet CSMA/CD, WiFi CSMA/CA) to coordinate access and reduce collisions.
    \item \textbf{MAC Addressing}: Frames contain physical (MAC) addresses which are distinct from IP addresses and identify specific network interfaces.
\end{itemize}

\subsubsection{Reliable Delivery}

\begin{itemize}
    \item \textbf{Definition}: Ensuring a datagram is delivered across a link without errors.
    \item \textbf{Mechanism}: Utilizes ARQ (ACKs and retransmissions).
    \item \textbf{Graduate-Level Trade-Off}:
          Even though TCP provides end-to-end reliability, local per-link reliability is useful on \emph{high-error links} (e.g., wireless). Local retransmission avoids expensive global retransmission.
    \item \textbf{Optimization}: Links with very low bit-error rates (e.g., fiber) typically omit reliable delivery to avoid unnecessary overhead.
\end{itemize}

\subsubsection{Error Detection and Correction}

\begin{itemize}
    \item \textbf{Definition}: Detecting or correcting bit-level corruption.
    \item \textbf{Mechanisms}:
          \begin{itemize}
              \item \emph{Error Detection}: CRC, checksums in the trailer.
              \item \emph{Error Correction}: Forward Error Correction (FEC) enables correction without retransmission.
          \end{itemize}
    \item \textbf{Implementation}: Usually hardware-accelerated for speed.
\end{itemize}

% ----------------------------------------------------------------------
\subsection{Implementation of the Link Layer}

\subsubsection{Location: The Network Adapter (NIC)}

Most link-layer functionality is implemented in specialized hardware known as the \textbf{Network Interface Controller (NIC)}.

\paragraph{Host Architecture Overview.}

\begin{verbatim}
      [ Host Architecture ]
               |
[ CPU / Memory / OS (Software) ]
      ^        |
      | (Bus)  |
      v        v
[ Network Adapter (NIC) ]
      |  - Link-layer controller chip
      |  - Local RAM buffer
      |  - DSP / modem components
      |
      v
[ Physical Link ]
\end{verbatim}

\subsubsection{Hardware–Software Interface}

\begin{itemize}
    \item \textbf{Hardware (Controller)}:
          Performs framing, error detection, MAC logic, and timing-sensitive operations at line rate (10~Gbps--100~Gbps).
    \item \textbf{Software (Driver)}:
          \begin{enumerate}
              \item Prepares link-layer addressing information.
              \item Activates and manages NIC hardware.
              \item Handles interrupts to pass received datagrams to upper layers.
          \end{enumerate}
\end{itemize}

\subsubsection{Data Flow}

\paragraph{Sending Side.}
\begin{enumerate}
    \item OS passes the datagram to the NIC.
    \item NIC encapsulates it in a frame, computes error-check bits.
    \item Bits are transmitted onto the physical medium.
\end{enumerate}

\paragraph{Receiving Side.}
\begin{enumerate}
    \item NIC receives raw bits and reconstructs the frame.
    \item Performs error detection and discards or accepts the frame.
    \item If valid, NIC extracts the datagram and interrupts the OS to pass it upward.
\end{enumerate}

% ----------------------------------------------------------------------
\subsection{Summary and Reinforcement}

\paragraph{Key Terms.}

\begin{itemize}
    \item \textbf{Frame}: Link-layer protocol data unit.
    \item \textbf{Node}: Device implementing link-layer logic.
    \item \textbf{MAC Protocol}: Rules for shared-medium access.
    \item \textbf{NIC}: Hardware interface implementing link-layer operations.
\end{itemize}

\paragraph{Core Relationships.}

\begin{itemize}
    \item \textbf{Link vs.\ Network Layer}: Link layer = node-to-node; network layer = end-to-end.
    \item \textbf{Software vs.\ Hardware}: Link-layer boundary marks the shift from OS software to NIC hardware.
\end{itemize}

\paragraph{Key Takeaways.}

\begin{itemize}
    \item Link-layer protocols are independent across different links.
    \item Services vary: wireless often requires reliability; fiber typically does not.
    \item High-speed hardware implementations (ASICs, FPGAs) are required for error checking and MAC functions.
\end{itemize}

\section{Error-Detection and -Correction Techniques (Section 6.2)}

\subsection{The Fundamental Problem: Bit Errors}

\subsubsection{Concept Overview}

At the physical layer, data is transmitted as analog signals that are vulnerable to
\textbf{attenuation}, \textbf{thermal noise}, and \textbf{electromagnetic interference (EMI)}.
These disturbances may cause the receiver to misinterpret a transmitted bit, such as reading a
``0'' as a ``1''.

The Link Layer must therefore provide mechanisms to \emph{detect} and, in some cases,
\emph{correct} these bit errors locally, rather than forcing upper layers such as IP or TCP
to repeatedly handle corrupted payloads. This is achieved by introducing \textbf{redundancy}
into the transmitted frame.

\subsubsection{The Error-Detection Model}

Let:
\begin{itemize}
    \item $D$: Data protected by error checking, may include header fields
    \item $EDC$: Error-detection/-correction bits (redundancy).
\end{itemize}

The sender transmits the bit pattern:
\[
    [\, D \ | \ EDC \,].
\]

The receiver obtains:
\[
    [\, D' \ | \ EDC' \,],
\]
where $D'$ and $EDC'$ may differ from $D$ and $EDC$ due to bit flips.

The receiver computes a function $f(D')$ and compares it with $EDC'$.

\begin{itemize}
    \item If they match: accept the frame.
    \item If they differ: an error is detected.
\end{itemize}

\paragraph{Error-Detection Scenario (Conceptual Diagram)}

\begin{verbatim}
      Sender                                       Receiver
   Datagram D                                 Datagram D'
        |                                          ^
        v                                          |
   [Compute EDC]                           [Check for Errors]
        |                                          ^
        v                                          |
   [Frame: D | EDC] --(Bit Errors on Link)--> [Frame: D' | EDC']
\end{verbatim}

\subsubsection{Key Limitation: Undetected Errors}

No error-detection method is perfect. There is always a small probability that
bit flips occur in such a way that the corrupted data $D'$ produces a check value
matching the corrupted $EDC'$. Such cases are called \textbf{undetected errors}.

% ---------------------------------------------------------------------
\subsection{Parity Checks}

\subsubsection{Single-Bit Parity}

\paragraph{Mechanism.}
\begin{itemize}
    \item \textbf{Even parity}: append a bit so total \# of 1s is even.
    \item \textbf{Odd parity}: append a bit so total \# of 1s is odd.
\end{itemize}

\paragraph{Detection Capability.}
\begin{itemize}
    \item Detects \textbf{any odd number} of bit flips.
    \item Fails for \textbf{even-numbered} bit flips (e.g., 2 flipped bits).
\end{itemize}

\paragraph{Intuition.}
Simple and cheap but weak, especially for \emph{burst errors} where multiple adjacent bits corrupt.

\subsubsection{Two-Dimensional Parity}

\paragraph{Mechanism.}
\begin{itemize}
    \item Organize data $D$ into an $i \times j$ bit matrix.
    \item Compute parity for each row and each column.
    \item Add a possible overall parity bit.
\end{itemize}

\paragraph{Correction Capability.}
A single flipped bit causes a parity error in:
\begin{itemize}
    \item its \textbf{row}, and
    \item its \textbf{column}.
\end{itemize}

The intersection identifies the erroneous bit, enabling
\textbf{Forward Error Correction (FEC)} by flipping it back.

\paragraph{Two-Dimensional Parity Illustration}

\begin{verbatim}
          Data Columns         Row Parity
       +-------------------+
Row 1: | 1 0 1 0 1 | 1 |
Row 2: | 1 0 1 1 1 | 0 |
Row 3: | 0 1 1 1 0 | 1 |
       +-------------------+
Col    | 1 0 1 1 1 | 0 |
Parity
\end{verbatim}

If, for example, the bit at (Row 2, Col 2) flips from 1 to 0, both Row 2 and Col 2 show parity errors.
\textbf{Intersection = correctable bit position}.

% ---------------------------------------------------------------------
\subsection{Checksumming Methods}

\subsubsection{Concept Overview}

Checksums are lightweight, software-friendly error-detection methods predominantly used in the
Network (IP) and Transport (TCP/UDP) layers, where performance must be balanced with acceptable detection strength.

\subsubsection{The Internet Checksum (RFC 1071)}

\paragraph{Sender Algorithm.}
\begin{enumerate}
    \item Treat the segment as a sequence of 16-bit integers.
    \item Sum them using \textbf{1's complement arithmetic} (wraparound carry).
    \item Take the 1's complement of the final sum to form the checksum.
\end{enumerate}

\paragraph{Receiver Algorithm.}
\begin{enumerate}
    \item Compute the 1's complement sum of all 16-bit words, including the checksum.
    \item If the result is all 1s (i.e., $1111\ldots1111$), no errors are detected.
    \item Otherwise, an error is detected.
\end{enumerate}

\paragraph{Intuition.}
1's complement arithmetic allows uniform behavior across different byte-order conventions.

% ---------------------------------------------------------------------
\subsection{Cyclic Redundancy Check (CRC)}

\subsubsection{Concept Overview}

CRC is the dominant error-detection technique in the Link Layer (Ethernet, WiFi).
It relies on \textbf{polynomial arithmetic modulo 2}.
Hardware implementations (e.g., linear feedback shift registers) make CRC extremely fast and robust.

\subsubsection{Mathematical Framework}

\begin{itemize}
    \item Treat data $D$ as a polynomial $D(x)$ over GF(2).
    \item Let $G(x)$ (generator polynomial) be a publicly defined $(r+1)$-bit pattern.
    \item Goal: append an $r$-bit residue $R$ so that:
\end{itemize}

\[
    \frac{D \cdot 2^r \oplus R}{G} = n, \qquad \text{remainder } 0.
\]

\subsubsection{Sender Algorithm}

\begin{enumerate}
    \item Append $r$ zeros to $D$ to form $D \cdot 2^r$.
    \item Divide ($D \cdot 2^r$) by $G$ using modulo-2 division.
    \item The remainder is the CRC value $R$.
    \item Transmit $[\, D \ | \ R \,]$.
\end{enumerate}

\subsubsection{Receiver Algorithm}

\begin{enumerate}
    \item Receive $[\, D' \ | \ R' \,]$.
    \item Divide the whole bitstring by $G$.
    \item If the remainder is zero, frame is assumed uncorrupted; otherwise, error detected.
\end{enumerate}

\subsubsection{Example: Polynomial Division}

Let $D = 101110$ and $G = 1001$ (so $r=3$).

\paragraph{Step 1:} Append 3 zeros:
\[
    101110000.
\]

\paragraph{Step 2:} Perform modulo-2 division:

\begin{verbatim}
              101011   <-- Quotient (unused)
         ___________
 1001 ) 101110000
        1001
        ----
         0101
         0000
         ----
          1010
          1001
          ----
           0110
           0000
           ----
            1100
            1001
            ----
             1010
             1001
             ----
              011   <-- Remainder R
\end{verbatim}

\paragraph{Result.}
Sender transmits:
\[
    [\,101110\ 011\,].
\]

\subsubsection{Practical Power of CRC}

CRC-32 (Ethernet) detects:
\begin{itemize}
    \item All burst errors of length $< r+1$.
    \item Any odd number of bit errors.
\end{itemize}

Probability of undetected error is typically $<10^{-12}$ for real-world frames.

% ---------------------------------------------------------------------
\subsection{Summary and Reinforcement}

\subsubsection{Key Terms}

\begin{itemize}
    \item \textbf{BER (Bit Error Ratio)}: Probability a bit is received incorrectly.
    \item \textbf{FEC (Forward Error Correction)}: Receiver-based correction without retransmission.
    \item \textbf{Modulo-2 Arithmetic}: XOR-based arithmetic with no carries.
    \item \textbf{Generator Polynomial $G(x)$}: Divisor used for CRC.
\end{itemize}

\subsubsection{Core Relationships}

\begin{itemize}
    \item Parity: simplest but weakest.
    \item Checksums: software-efficient but moderate detection capability.
    \item CRC: hardware-efficient, extremely strong detection.
\end{itemize}

\subsubsection{Key Insights}

\begin{itemize}
    \item Reliability is probabilistic; no method guarantees perfect detection.
    \item FEC reduces latency on noisy or high-latency links (e.g., wireless, satellite).
    \item CRC hardware uses XOR gates and LFSRs for high-speed computation.
\end{itemize}


\section{Multiple Access Links and Protocols (Section 6.3)}

\subsection{The Multiple Access Problem}

\subsubsection{Concept Overview}

Up to this point, we have examined \emph{point-to-point} links (e.g., a dedicated wire between two routers).
However, many networks rely on a \textbf{broadcast link}, where \emph{multiple nodes share a common communication medium}. Examples include:
\begin{itemize}
    \item Traditional shared-medium Ethernet,
    \item WiFi (802.11),
    \item Satellite networks,
    \item Shared coaxial cable systems.
\end{itemize}

\paragraph{Cocktail Party Analogy.}
If multiple people in a room speak at the same time, no one is understood. Similarly, nodes on a broadcast channel must coordinate \emph{who transmits when}.

\subsubsection{Collisions}

When two or more nodes transmit simultaneously, their signals interfere, producing a \textbf{collision}.
\begin{itemize}
    \item Result: the signal becomes garbled and the frame is lost.
    \item Consequence: channel bandwidth is wasted during corrupted transmission.
\end{itemize}

\subsubsection{The Multiple Access Protocol}

A \textbf{Multiple Access Protocol} is a distributed algorithm that determines how nodes share a broadcast channel.

\paragraph{Ideal Characteristics.}
For a channel of rate $R$ with $M$ nodes, a good protocol should:
\begin{enumerate}
    \item Provide throughput $R$ when one node transmits and $R/M$ per node at full load,
    \item Be decentralized (no master node),
    \item Be simple and inexpensive to implement.
\end{enumerate}

% ---------------------------------------------------------------------
\subsection{Taxonomy of Multiple Access Protocols}

Multiple access protocols fall into three broad families:

\begin{enumerate}
    \item \textbf{Channel Partitioning Protocols}:
          \begin{enumerate}
              \item divide channel into sections (time/frequency) and assign exclusively.
              \item allocate piece to node for exclusive use
          \end{enumerate}
    \item \textbf{Random Access Protocols}: allow nodes to transmit at will; detect/recover from collisions.
          \begin{enumerate}
              \item allow nodes to transmit at will; detect/recover from collisions.
          \end{enumerate}
    \item \textbf{Taking-Turns Protocols}:
          \begin{enumerate}
              \item nodes transmit in a coordinated sequence.
          \end{enumerate}
\end{enumerate}

% ---------------------------------------------------------------------
\subsection{Channel Partitioning Protocols}

\subsubsection{TDMA (Time Division Multiple Access)}

\paragraph{Mechanism.}
\begin{itemize}
    \item Channel divided into frames with $N$ fixed time slots.
    \item Each node transmits only in its assigned slot.
    \item Slots that're unused go idle
\end{itemize}

\paragraph{Analysis.}
\begin{itemize}
    \item \textbf{Pro}: No collisions; perfect fairness.
    \item \textbf{Con}: Highly inefficient at low load—many slots may be idle.
\end{itemize}

\subsubsection{FDMA (Frequency Division Multiple Access)}

\paragraph{Mechanism.}
\begin{itemize}
    \item The total bandwidth $R$ is divided into $N$ frequency bands of width $R/N$.
    \item Each node transmits on its dedicated band.
    \item Unused transmission time in frequency bands go idle.
\end{itemize}

\paragraph{Analysis.}
Similar pros and cons as TDMA: high efficiency at heavy load but wasteful at light load.

% ---------------------------------------------------------------------
\subsection{Random Access Protocols}

Optimized for \textbf{bursty} traffic. Nodes transmit at the full channel rate $R$.
If a collision occurs, the node waits for a random delay and retransmits.

% -----------------------------
\subsubsection{Slotted ALOHA}

\paragraph{Mechanism.}
\begin{itemize}
    \item Time is divided into fixed-size slots of duration $L/R$.
    \item Nodes are synchronized to slot boundaries.
    \item A node transmits at the \emph{beginning} of the next slot when it has a frame.
    \item On collision, nodes retransmit with probability $p$ in future slots.
\end{itemize}

\paragraph{Efficiency Analysis.}
Let each of $N$ nodes attempt transmission with probability $p$:
\[
    \Pr(\text{Node A succeeds}) = p(1-p)^{N-1},
\]
\[
    \Pr(\text{Any node succeeds}) = Np(1-p)^{N-1}.
\]

Maximizing with respect to $p$ gives:
\[
    \text{Maximum Efficiency} = \frac{1}{e} \approx 0.37.
\]

\paragraph{Intuition.}
About 37\% of slots carry a successful transmission; the rest are idle or suffer collisions.

% -----------------------------
\subsubsection{Pure ALOHA (Unslotted)}

\paragraph{Mechanism.}
Nodes transmit immediately when a frame is ready—no slot synchronization.

\paragraph{Collision Window.}
A transmission at time $t_0$ collides with any transmission in $(t_0 - 1,\, t_0 + 1)$.

\paragraph{Efficiency.}
\[
    \text{Max Efficiency} = \frac{1}{2e} \approx 0.18.
\]

% -----------------------------
\subsubsection{CSMA (Carrier Sense Multiple Access)}

\paragraph{Philosophy.}
``Listen before transmitting.''

\paragraph{Mechanism.}
\begin{enumerate}
    \item If the channel is idle, transmit.
    \item If the channel is busy, wait until idle.
\end{enumerate}

\paragraph{Why Collisions Still Occur?}
\textbf{Propagation delay}.
A node may sense the channel idle because another node’s signal has not yet arrived.

% -----------------------------
\subsubsection{CSMA/CD (Carrier Sense Multiple Access with Collision Detection)}

\paragraph{Philosophy.}
``Listen while transmitting.''

\paragraph{Mechanism.}
\begin{enumerate}
    \item If channel idle → transmit.
    \item During transmission, monitor channel energy.
    \item If a collision is detected:
          \begin{itemize}
              \item Abort transmission.
              \item Send a 48-bit jam signal.
          \end{itemize}
    \item Use \textbf{binary exponential backoff}:
          \[
              K \in \{0,1,\dots,2^n - 1\},
          \]
          where $n$ is the number of collisions so far. Wait $K \cdot 512$ bit-times.
\end{enumerate}

\paragraph{Efficiency.}
\[
    \text{Efficiency} = \frac{1}{1 + 5(d_{prop}/d_{trans})}.
\]

\paragraph{Insight.}
Efficiency increases when:
\begin{itemize}
    \item propagation delay $\to 0$,
    \item frame size $\to \infty$.
\end{itemize}

Thus CSMA/CD is most effective on short links with large frames.

% ---------------------------------------------------------------------
\subsection{Taking-Turns Protocols}

\subsubsection{Polling}

\paragraph{Mechanism.}
A master node polls each slave node in turn:
``Do you have data to send?''

\paragraph{Trade-offs.}
\begin{itemize}
    \item No collisions.
    \item Polling overhead and single point of failure.
\end{itemize}

\subsubsection{Token Passing}

\paragraph{Mechanism.}
A token frame circulates among nodes.
A node may transmit only when holding the token.

\paragraph{Trade-offs.}
\begin{itemize}
    \item Highly efficient at high load.
    \item Token recovery is complex if a node crashes while holding it.
\end{itemize}

% ---------------------------------------------------------------------
\subsection{DOCSIS: Multiple Access for Cable Internet}

DOCSIS divides communication into downstream and upstream paths:

\paragraph{Downstream (CMTS → Cable Modems).}
Broadcast channel—no contention.

\paragraph{Upstream (Cable Modems → CMTS).}
Multiple modems contend for access:
\begin{itemize}
    \item Uses \textbf{TDMA} for actual data slots.
    \item Uses \textbf{Slotted ALOHA} for request messages.
\end{itemize}

\paragraph{Process.}
\begin{enumerate}
    \item Modems send MAP requests via Slotted ALOHA.
    \item The CMTS assigns specific TDMA slots for collision-free data transmission.
\end{enumerate}

% ---------------------------------------------------------------------
\subsection{Summary and Reinforcement}

\subsubsection{Key Terms}

\begin{itemize}
    \item \textbf{Collision}: simultaneous transmissions causing interference.
    \item \textbf{Carrier Sense}: checking whether the channel is idle.
    \item \textbf{Exponential Backoff}: dynamic waiting time based on congestion.
    \item \textbf{Propagation Delay}: physical travel time of a signal.
\end{itemize}

\subsubsection{Core Relationships}

\begin{itemize}
    \item \textbf{ALOHA}: efficient at low load, collapses at high load.
    \item \textbf{TDMA/FDMA}: inefficient at low load, ideal at high load.
    \item \textbf{CSMA/CD}: robust across medium–high loads.
\end{itemize}

\subsubsection{Key Insights}

\begin{itemize}
    \item No protocol is universally best—performance depends on traffic patterns.
    \item Longer physical links reduce CSMA/CD efficiency (delayed collision detection).
    \item Exponential backoff is essential to avoid congestion collapse in random-access systems.
\end{itemize}


\section{Switched Local Area Networks (Section 6.4)}

\subsection{Link-Layer Addressing and ARP (6.4.1)}

\subsubsection{MAC Addresses: The Physical Identity}

\paragraph{Concept Overview.}

IP addresses are hierarchical and location-dependent (like a postal address), whereas
\textbf{MAC addresses} are flat, permanent identifiers (like a fingerprint or Social Security Number).
A MAC address is burned into the NIC ROM and remains constant regardless of network location; a device moved from New York to Tokyo receives a new IP address but retains the same MAC address.

\paragraph{Technical Mechanism.}

\begin{itemize}
    \item \textbf{Format:} 48 bits (6 bytes), written in hexadecimal (e.g., \texttt{1A:2B:3C:4D:5E:6F}).
    \item \textbf{Allocation:} IEEE assigns Organizationally Unique Identifiers (OUIs) for the first 24 bits; vendors assign the remaining 24 bits.
    \item \textbf{Broadcast Address:} \texttt{FF:FF:FF:FF:FF:FF}. Frames sent to this address are processed by \emph{all} devices on the LAN.
\end{itemize}

\paragraph{Why Two Addresses (IP and MAC)?}

\begin{itemize}
    \item \textbf{Layering Abstraction:}
          Network-layer routing must be independent of the underlying link technology.
          Using MAC-only addressing would require global flat routing tables of size proportional to the number of hosts.
    \item \textbf{Modularity:}
          Using IP addresses exclusively at the link layer would require every link protocol to understand IP addressing.
\end{itemize}

% ---------------------------------------------------------------------
\subsubsection{ARP: Address Resolution Protocol}

\paragraph{The Problem.}

A sender knows the destination's \emph{IP address} but requires the \emph{MAC address} to send a frame on a LAN.
ARP solves:
\[
    192.168.1.5 \longrightarrow 00:13:49:AE:F2:10.
\]

\paragraph{Mechanism: The ARP Protocol.}

\begin{itemize}
    \item \textbf{ARP Table:} A cache stored in RAM containing entries of the form:
          \[
              \langle \text{IP Address}, \text{MAC Address}, \text{TTL} \rangle.
          \]
          TTL ensures entries expire (commonly after 20 minutes).

    \item \textbf{ARP Query (Broadcast):}
          If an IP address is unknown, the host broadcasts an ARP request:
          \[
              \text{Dest MAC} = \texttt{FF:FF:FF:FF:FF:FF}.
          \]

    \item \textbf{ARP Response (Unicast):}
          The owner of the requested IP replies directly with its MAC address.

    \item \textbf{Caching:}
          Sender stores the discovered mapping in its ARP table.
\end{itemize}

\paragraph{ARP Operation Diagram.}

\begin{verbatim}
      [Node A]                                         [Node B]
 (IP: A.A.A.A, MAC: 11-11)                       (IP: B.B.B.B, MAC: 22-22)
       |                                                 |
       |  1. Check ARP Table (Miss)                      |
       |------------------------------------------------>|
       |  2. ARP Query (Broadcast Frame)                 |
       |     [Src: 11-11, Dst: FF-FF, TargetIP: B.B.B.B] |
       |                                                 |
       |                                                 | 3. B identifies its IP
       |<------------------------------------------------|
       |  4. ARP Response (Unicast Frame)                |
       |     [Src: 22-22, Dst: 11-11, MyMAC: 22-22]      |
       |                                                 |
       |  5. Update ARP Table                            |
       v
\end{verbatim}

\subsubsection{Sending a Packet Off-Subnet}

If Node A sends to a destination Node Z on a remote subnet, \emph{Node A never ARPs for Node Z's MAC}.
Instead, Node A ARPs for the \textbf{First-Hop Router's MAC}.

\begin{itemize}
    \item \textbf{Destination MAC:} Router’s interface MAC.
    \item \textbf{Destination IP:} Node Z’s ultimate IP address.
\end{itemize}

The router removes the link-layer header, examines the IP packet, forwards to the next hop, and constructs a new link-layer frame for the next link.

% ---------------------------------------------------------------------
\subsection{Ethernet (6.4.2)}

\subsubsection{Overview and History}

Ethernet became the dominant wired LAN technology, surpassing Token Ring and ATM due to:
\begin{enumerate}
    \item \textbf{Simplicity:} No tokens or complex control mechanisms.
    \item \textbf{Evolution:} Speed increases (10M $\rightarrow$ 100M $\rightarrow$ 1G $\rightarrow$ 10G $\rightarrow$ 400G) without changing the frame structure.
\end{enumerate}

\paragraph{Topology.}
\begin{itemize}
    \item \textbf{Legacy:} Bus topology, coaxial cable, collisions common.
    \item \textbf{Modern:} Star topology, switches, full-duplex, no collisions.
\end{itemize}

\subsubsection{Ethernet Frame Structure}

\begin{verbatim}
+----------+----------+-----------+-----------+----------------+-------+
| Preamble | Dest MAC | Source MAC|   Type    | Data (Payload) |  CRC  |
| (8 Bytes)| (6 Bytes)| (6 Bytes) | (2 Bytes) | (46-1500 Bytes)|(4 B)  |
+----------+----------+-----------+-----------+----------------+-------+
\end{verbatim}

\paragraph{Fields.}

\begin{itemize}
    \item \textbf{Preamble (8 bytes):} Seven bytes of 10101010 and one byte of 10101011 for receiver clock synchronization.
    \item \textbf{Addresses:} Destination and source MAC addresses.
    \item \textbf{Type:} Identifies higher-layer protocol (e.g., IPv4, IPv6, ARP).
    \item \textbf{Data:} MTU = 1500 bytes; minimum 46 bytes (padding added if needed).
    \item \textbf{CRC:} 32-bit error detection. Frames failing CRC are dropped (Ethernet does not retransmit).
\end{itemize}

\subsubsection{Service Model}

\begin{itemize}
    \item \textbf{Connectionless:} No handshake.
    \item \textbf{Unreliable:} No ACKs; corrupted frames are silently discarded.
\end{itemize}

% ---------------------------------------------------------------------
\subsection{Link-Layer Switches (6.4.3)}

\subsubsection{Definition and Purpose}

A switch is a transparent Layer-2 packet switch that forwards frames based on MAC addresses.
Unlike hubs, switches isolate collision domains and support parallel transmissions.

\subsubsection{Switch Table}

Each switch maintains a table of entries:
\[
    \langle \text{MAC Address}, \text{Interface}, \text{Timestamp} \rangle.
\]

\subsubsection{Self-Learning Algorithm (Plug-and-Play)}

\begin{enumerate}
    \item \textbf{Frame Arrives:} On interface $x$ with source $S$ and destination $D$.
    \item \textbf{Learn:} Update: “MAC $S$ is on interface $x$”.
    \item \textbf{Forwarding:}
          \begin{itemize}
              \item If $D$ known and mapped to interface $y$:
                    \begin{itemize}
                        \item If $x = y$: drop (same segment).
                        \item If $x \neq y$: forward to $y$ only.
                    \end{itemize}
              \item If $D$ unknown: \textbf{flood} to all interfaces except $x$.
          \end{itemize}
\end{enumerate}

\paragraph{Self-Learning Sequence.}

\begin{verbatim}
Initial State: Empty Table []

Event 1: A (on Port 1) sends to B (on Port 2).
   -> Switch Learns: {A: Port 1}
   -> Switch Unknown Dest B: Floods to Ports 2, 3, 4.

Event 2: B replies to A.
   -> Switch Learns: {B: Port 2}
   -> Switch Checks Dest A: Known! (Port 1)
   -> Switch Forwards ONLY to Port 1.
\end{verbatim}

\subsubsection{Switches vs. Routers}

\begin{center}
    \begin{tabular}{|l|c|c|c|}
        \hline
        \textbf{Feature}  & \textbf{Hub}     & \textbf{Switch (L2)}    & \textbf{Router (L3)}          \\
        \hline
        Traffic Isolation & No               & Yes (collision domains) & Yes (broadcast domains)       \\
        Addressing        & None             & MAC (flat)              & IP (hierarchical)             \\
        Routing           & None             & Self-learning           & Routing protocols (OSPF, BGP) \\
        Processing        & Bit-level repeat & Frame inspection        & Packet inspection + TTL       \\
        Throughput        & Shared           & High (parallel)         & Lower (complex processing)    \\
        \hline
    \end{tabular}
\end{center}

\paragraph{Key Insight.}
Switches scale well and are plug-and-play, but routers are essential to prevent broadcast storms and implement hierarchical routing.

% ---------------------------------------------------------------------
\subsection{Virtual Local Area Networks (VLANs) (6.4.4)}

\subsubsection{Motivation}

In a traditional LAN:
\begin{itemize}
    \item Excessive broadcast traffic wastes CPU cycles,
    \item Reorganization requires physical rewiring,
    \item All users can see broadcast frames (security issue).
\end{itemize}

\subsubsection{Port-Based VLANs}

A VLAN-capable switch assigns ports to logical groups.

\begin{itemize}
    \item Ports 1--8 $\rightarrow$ VLAN 10 (Engineering).
    \item Ports 9--16 $\rightarrow$ VLAN 20 (Sales).
\end{itemize}

Broadcasts remain within VLAN boundaries.
Routing is required for communication across VLANs.

\subsubsection{VLAN Trunking (802.1Q)}

When VLANs span multiple switches, a trunk link carries frames for all VLANs using the \textbf{802.1Q tag}.

\paragraph{802.1Q Frame Tag.}

\begin{itemize}
    \item \textbf{TPID:} \texttt{0x8100} (identifies tagged frames).
    \item \textbf{TCI:} Contains 12-bit VLAN ID.
\end{itemize}

\paragraph{Operation:}
\begin{enumerate}
    \item Switch A receives an untagged frame from VLAN 10.
    \item Switch A inserts 802.1Q tag and forwards over trunk.
    \item Switch B reads tag (ID=10), strips it, and forwards to VLAN 10 ports.
\end{enumerate}

\paragraph{802.1Q Tag Illustration.}

\begin{verbatim}
Original: [Dest][Src][Type][Data][CRC]
Tagged:   [Dest][Src][Type 8100][VLAN ID][Type][Data][CRC]
                     ^-- Added 4 Bytes --^
\end{verbatim}

% ---------------------------------------------------------------------
\subsection{Summary and Reinforcement}

\subsubsection{Key Terms}

\begin{itemize}
    \item \textbf{MAC Address}: 48-bit hardware identifier.
    \item \textbf{ARP}: Local broadcast-based IP $\rightarrow$ MAC resolution.
    \item \textbf{Switching Table}: $\langle$MAC, Interface, TTL$\rangle$ mapping.
    \item \textbf{Collision Domain}: Range where collisions may occur; broken by switches.
    \item \textbf{Broadcast Domain}: Range where broadcast frames propagate; broken by routers.
    \item \textbf{VLAN}: Logical segmentation of Layer-2 networks.
\end{itemize}

\subsubsection{Core Relationships}

\begin{itemize}
    \item IP gets packets to the correct subnet; MAC gets packets to the correct interface.
    \item Switches isolate collisions; hubs do not.
    \item VLANs separate groups logically; routers interconnect them.
\end{itemize}

\subsubsection{Key Insights}

\begin{itemize}
    \item \textbf{Self-learning} allowed Ethernet to scale without manual configuration.
    \item \textbf{Plug-and-play:} Ethernet works instantly compared to IP, which requires configuration (DHCP/manual).
    \item \textbf{Scalability Limit:} Broadcast traffic grows with $N$; routers are required to partition large L2 domains.
\end{itemize}


\section{Link Virtualization: A Network as a Link Layer (Section 6.5)}

\subsection{The Evolution of the ``Link'' Abstraction}

\subsubsection{Concept Overview}

Throughout this course, our definition of a ``link'' has expanded from a simple physical connection to a sophisticated logical abstraction. This evolution culminates in the concept of \textbf{Link Virtualization}.

\begin{itemize}
    \item \textbf{Phase 1: The Physical Wire} – A copper wire directly connecting two hosts.
    \item \textbf{Phase 2: The Shared Medium} – A medium shared by many nodes (Ethernet bus, WiFi spectrum) requiring multiple access protocols.
    \item \textbf{Phase 3: The Switched Fabric} – A switched LAN acting as a transparent fabric connecting hosts.
    \item \textbf{Phase 4: The Network as a Link (Virtualization)} – An entire network (e.g., the PSTN or MPLS core) can act as a \emph{single logical link} to the upper layers.
\end{itemize}

\subsubsection{The Dial-Up Analogy}

A home user connecting to an ISP via dial-up experiences:

\begin{itemize}
    \item \textbf{Physical Reality:} The signal traverses the PSTN (switches, copper, fiber, multiplexers).
    \item \textbf{Link Layer View:} The operating system sees a simple point-to-point link to the ISP.
    \item \textbf{Virtualization:} The PSTN \emph{virtualizes} a link for the Internet.
\end{itemize}

% ---------------------------------------------------------------------
\subsection{Multiprotocol Label Switching (MPLS)}

\subsubsection{Definition and Purpose}

Multiprotocol Label Switching (MPLS) is a forwarding technology using short, fixed-length \textbf{labels} rather than destination IP address lookups.

\begin{itemize}
    \item \textbf{Original Goal:} Speed up packet forwarding (IP lookup was slow in 1990s hardware).
    \item \textbf{Modern Use:} Traffic Engineering (TE), Virtual Private Networks (VPNs), Fast Reroute.
    \item \textbf{Hybrid Nature:} MPLS sits between Layers 2 and 3 (``Layer 2.5'').
\end{itemize}

\subsubsection{The MPLS Header}

A new 32-bit MPLS header is inserted between the link-layer header (e.g., Ethernet) and the IP header.

\begin{verbatim}
+------------------+-------------+-----------+----------------------+
| Layer 2 Header   | MPLS Header | IP Header |      Payload         |
| (e.g., Ethernet) | (32 bits)   |           |                      |
+------------------+-------------+-----------+----------------------+
\end{verbatim}

\paragraph{Header Fields (32 bits):}

\begin{enumerate}
    \item \textbf{Label (20 bits)} – Used to forward packets.
    \item \textbf{Exp (3 bits)} – Typically used for QoS / Class of Service.
    \item \textbf{S (1 bit)} – Bottom-of-stack indicator (MPLS supports label stacking).
    \item \textbf{TTL (8 bits)} – Prevents forwarding loops (same function as IP TTL).
\end{enumerate}

\subsubsection{Mechanism: Label Switched Routers (LSRs)}

Routers supporting MPLS are called \textbf{Label Switched Routers}. They forward based solely on the label rather than inspecting destination IP addresses.

\paragraph{MPLS Forwarding Table:}

Traditional IP forwarding table:

\[
    \text{Dest IP Prefix} \rightarrow \text{Out Interface}.
\]

MPLS table:

\[
    [\text{In Label}, \text{In Interface}] \rightarrow [\text{Out Label}, \text{Out Interface}].
\]

\paragraph{Forwarding Process:}

\begin{enumerate}
    \item \textbf{Ingress (Label Push):}
          The ingress router assigns a label (e.g., 10) based on the IP header and attaches the MPLS header.
    \item \textbf{Transit (Label Swap):}
          A transit LSR swaps the incoming label with the appropriate outgoing label.
    \item \textbf{Egress (Label Pop):}
          The final router strips the label and forwards the remaining IP packet normally.
\end{enumerate}

\paragraph{MPLS Forwarding Example.}

\begin{verbatim}
     [R1: Ingress]          [R2: Transit]          [R3: Egress]
      (IP Router)            (LSR)                  (IP Router)
          |                      |                      |
[IP Packet] ---> [ PUSH Label 6 ] ---> [ SWAP 6->12 ] ---> [ POP Label ] -> [IP Packet]
          |       [ Label: 6 ]   |      [ Label: 12 ]   |
          |                      |                      |
     [Lookup Dest IP]       [Lookup Label 6]       [Lookup Dest IP]
\end{verbatim}

% ---------------------------------------------------------------------
\subsection{Why MPLS? (The Value Proposition)}

If routers can now forward IP packets at wire speed, why use MPLS?

\subsubsection{Traffic Engineering}

\begin{itemize}
    \item \textbf{IP Limitation:}
          IP routing (OSPF/BGP) selects the \emph{shortest path}, even if congested.
    \item \textbf{MPLS Solution:}
          Operators define explicit traffic paths:
          \begin{itemize}
              \item Voice traffic (Label 100) via fast path.
              \item Bulk email (Label 200) via slower, uncongested path.
          \end{itemize}
    \item \textbf{Mechanism:}
          Establish a Label Switched Path (LSP) that behaves like a virtual circuit.
\end{itemize}

\subsubsection{Fast Reroute (Restoration)}

\begin{itemize}
    \item MPLS can precompute backups for every path.
    \item Failover occurs in under 50ms.
    \item OSPF convergence may take several seconds.
\end{itemize}

\subsubsection{Virtual Private Networks (VPNs)}

\begin{itemize}
    \item ISPs use MPLS to create traffic-isolated logical networks for customers.
    \item Customer A’s labeled packets and Customer B’s labeled packets share the same physical infrastructure but remain isolated.
    \item Each customer views the ISP backbone as a virtual private link between branch offices.
\end{itemize}

% ---------------------------------------------------------------------
\subsection{Summary and Reinforcement}

\subsubsection{Key Terms and Definitions}

\begin{itemize}
    \item \textbf{MPLS (Multiprotocol Label Switching):}
          Forwarding technology using short labels rather than IP prefixes.
    \item \textbf{LSR (Label Switched Router):}
          MPLS-capable router performing label swapping.
    \item \textbf{LSP (Label Switched Path):}
          Predefined forwarding path taken by MPLS labeled packets.
    \item \textbf{Traffic Engineering:}
          Steering traffic along controlled paths to optimize performance.
\end{itemize}

\subsubsection{Core Relationships}

\begin{itemize}
    \item MPLS is often described as \textbf{Layer 2.5}: encapsulated by L2, encapsulating L3.
    \item IP is simple and stateless but limited; MPLS is more complex but offers flexible routing, TE, VPNs.
\end{itemize}

\subsubsection{Key Insights / Takeaways}

\begin{itemize}
    \item \textbf{Decoupling Routing from Forwarding:}
          MPLS separates \emph{path selection} (at ingress) from \emph{hop-by-hop forwarding} (label swapping).
    \item \textbf{``Link'' as a Construct:}
          MPLS demonstrates that a “link” to IP may represent an entire network of LSRs and forwarding rules.
    \item \textbf{SDN Connection:}
          Software-Defined Networking generalizes the MPLS concept: forwarding entries are programmed by a controller rather than by distributed signaling (LDP/RSVP-TE).
\end{itemize}

\section{Data Center Networking (Section 6.6)}

\subsection{Concept Overview}

A \textbf{Data Center Network (DCN)} is the internal, complex computer network that interconnects the tens to hundreds of thousands of computing devices (hosts) within a massive data center \cite{13885}. DCNs are the backbone of modern cloud computing and large Internet service platforms like search, streaming video, e-commerce, social media, and large-scale machine learning systems \cite{13886}. Modern DCNs must simultaneously support high throughput, low latency, failure resilience, and massive scalability.

\paragraph*{Purpose and Applications.}
DCNs support several large-scale services:
\begin{itemize}
    \item \cite{13886} \textbf{Content Provision:} Serving Web pages, search results, e-mail, and streaming video to external clients.
    \item \cite{13886} \textbf{Massively Parallel Computing:} Supporting distributed computing tasks such as search index construction, data analytics, and machine learning pipelines.
    \item \cite{13886} \textbf{Cloud Computing/Services:} Providing IaaS, PaaS, and SaaS for platforms such as AWS, Azure, and Google Cloud.
\end{itemize}

\paragraph*{Underlying Mechanisms (Hosts and Racks).}
\begin{itemize}
    \item \cite{13886} The fundamental work unit in a DCN is the \textbf{host} (or blade), a commodity server with CPU, RAM, and local storage.
    \item \cite{13886} Hosts are physically organized into \textbf{racks}, typically containing 20--40 hosts.
    \item Each rack connects upward into the DCN via a \textbf{Top-of-Rack (TOR) switch}, which aggregates host traffic and forwards it toward higher-tier switching layers.
\end{itemize}

\paragraph*{Workload Characteristics (From the Textbook).}
According to Section 6.6, DCN workloads display the following important patterns:
\begin{itemize}
    \item \textbf{Highly parallel applications} generate \emph{shuffle traffic}, where many hosts exchange intermediate data.
    \item \textbf{Data analytics jobs} exhibit \emph{query fan-out}: a single user request may trigger hundreds of internal queries.
    \item \textbf{East-West traffic dominates}: Most DCN traffic flows between hosts \emph{within} a data center rather than between the data center and the external Internet.
\end{itemize}

These patterns heavily influence DCN architecture and routing strategies.

% ------------------------------------------------------------------

\subsection{Technical Mechanism: Data Center Architectures}

The internal network topology is predominantly hierarchical and designed for massive scale and high throughput \cite{13778}. This structure resembles a \textbf{Clos network} (also called a \textbf{fat-tree topology}), providing rich multi-path connectivity and high bisection bandwidth \cite{13778,13779}. Conceptually, a Clos fabric behaves like one giant non-blocking switch interconnecting all hosts.

\subsubsection{A. Hierarchical/Tiered Network Topology}

Traditional DCNs adopt a three-tier hierarchy beneath a set of border routers:

\begin{center}
    \begin{tabular}{|c|c|p{5cm}|p{4.5cm}|}
        \hline
        \textbf{Layer}       & \textbf{Element}     & \textbf{Function}                                & \textbf{Interconnection}                         \\
        \hline
        Edge                 & Hosts (Blades)       & Run applications and store data.                 & Connect to TOR switches (40--100 Gbps Ethernet). \\
        \hline
        Tier 3 (Access)      & TOR Switch           & Aggregates host traffic within a rack.           & Connects to Tier-2 switches.                     \\
        \hline
        Tier 2 (Aggregation) & Tier-2 Switches      & Aggregate traffic from dozens of TORs.           & Connect to Tier-1/core switches.                 \\
        \hline
        Tier 1 (Core)        & Core/Tier-1 Switches & Provide fabric-wide interconnection.             & Connect to Tier-2 switches and border routers.   \\
        \hline
        Border               & Border Routers       & Connect the DCN to ISPs and the public Internet. & Connect upward to external networks.             \\
        \hline
    \end{tabular}
\end{center}

\paragraph*{Oversubscription (Important Detail from Textbook).}
Traditional hierarchical DCNs often suffer from \textbf{oversubscription}, meaning:
\[
    \text{Capacity(host-to-host)} \; > \; \text{Capacity(available upward bandwidth)}.
\]
Typical oversubscription ratios in early DCNs were 5:1, 10:1, or even 20:1, severely limiting inter-rack throughput.

% ------------------------------------------------------------------

\subsubsection{B. The Challenge of Limited Host-to-Host Capacity}

\paragraph*{Problem.}
\cite{13778} The tiered architecture limits \textbf{bisection bandwidth}: the capacity supporting traffic between hosts in different halves of the network. This constraint is most severe when much of the traffic is \emph{cross-rack}, which is typical for distributed applications.

\paragraph*{Mechanism of Failure (Bottleneck).}
\cite{13778} A TOR switch may serve 40 hosts at 10 Gbps each (400 Gbps total), yet its uplink capacity to Tier 2 might be only 100 Gbps. When many hosts concurrently communicate externally or across racks, the uplink becomes an immediate bottleneck.

\paragraph*{Solution: Increased Interconnection \& Multi-Path Routing.}
\begin{itemize}
    \item Higher-tier switches are interconnected in mesh-like Clos/Fat-Tree patterns.
    \item \cite{13778,13779} Multi-path routing (especially via ECMP) distributes traffic across many equal-cost paths.
    \item \cite{13778} Increased redundancy simultaneously improves reliability and throughput.
\end{itemize}

\paragraph*{Clos/Fat-Tree Benefits (Explicit from Section 6.6).}
\begin{itemize}
    \item Uniform bandwidth between any two racks.
    \item Rich path diversity tolerates failures without service interruption.
    \item Scales naturally by adding more switching stages.
\end{itemize}

% ------------------------------------------------------------------

\subsubsection{C. Load Balancing}

\paragraph*{Mechanism.}
Incoming client requests targeting a public service IP address first reach a \textbf{Load Balancer} (typically a Layer 4 device).

\paragraph*{Function.}
\cite{13778} The load balancer distributes traffic across many internal hosts to balance load, reduce latency, and maximize utilization.

\paragraph*{Security and Addressing.}
\cite{13778} It also performs a NAT-like function, translating the public IP address into internal private host addresses, preventing direct external access to hosts.

% ------------------------------------------------------------------

\subsection{Practical and Intuitive View: Trends and Evolution}

\subsubsection{A. Centralized SDN Control and Management}

\paragraph*{Connection to Chapter 5.}
Unlike typical ISP networks (with per-router control), DCNs are owned and operated by a single organization. This allows for \textbf{centralized SDN control} \cite{13779,13778}.

\paragraph*{Mechanism.}
A logically centralized SDN controller programs flow tables for the DCN’s switches. This enables:
\begin{itemize}
    \item automated multi-path routing,
    \item dynamic load balancing,
    \item consistent network-wide security policy enforcement,
    \item fast failure recovery.
\end{itemize}

% ------------------------------------------------------------------

\subsubsection{B. Virtualization and Network Layer Abstraction}

\paragraph*{Problem: VM Mobility.}
VMs may move across racks for load balancing or hardware maintenance. Traditional IP/MAC addressing binds addresses to physical topology, complicating mobility.

\paragraph*{Mechanism: Flat Layer 2 Abstraction.}
To support seamless migration, the entire DCN is often treated as a \textbf{giant flat L2 network}. Instead of ARP floods:
\begin{itemize}
    \item a directory service maps VM IP $\rightarrow$ VM MAC $\rightarrow$ VM location,
    \item hosts query this directory for updated mappings.
\end{itemize}

This allows the DCN to behave like one massive Ethernet switch despite being physically multi-tiered.

% ------------------------------------------------------------------

\subsubsection{C. The Unique Challenge of Data Center Congestion Control}

\paragraph*{Physical Constraints.}
DCNs operate with extremely high bandwidth (40--100+ Gbps) and very low latency (microseconds) \cite{13780}.

\paragraph*{Congestion Problem.}
Classic TCP reacts only after packet loss, which is too slow at microsecond timescales. By the time loss occurs, queues may be overwhelmed.

\paragraph*{Custom Solutions (DCTCP/DCQCN).}
\cite{13780,13794} Modern DCNs use specialized transport protocols such as:
\begin{itemize}
    \item \textbf{DCTCP}: Uses ECN marking rather than loss to detect early congestion.
    \item \textbf{DCQCN}: A more advanced variant for RDMA over Converged Ethernet (RoCE).
\end{itemize}

These protocols reduce queue buildup, maintain extremely low latency, and improve throughput.

% ------------------------------------------------------------------

\subsection{Summary and Reinforcement}

\subsubsection{Key Terms and Definitions}

\begin{itemize}
    \item \cite{13778} \textbf{TOR Switch (Top-of-Rack):} The first-tier switch connecting hosts in a rack.
    \item \cite{13886} \textbf{Blade:} A thin, modular server unit; the fundamental computing element.
    \item \cite{13779} \textbf{Clos Network/Fat-Tree:} A multi-stage interconnection network providing high bisection bandwidth.
    \item \textbf{Bisection Bandwidth:} Maximum traffic capacity between any two halves of a network.
    \item \cite{13778} \textbf{Load Balancer:} Component distributing client requests across multiple servers while performing address translation.
\end{itemize}

\subsubsection{Core Relationships}

\begin{itemize}
    \item \cite{13778} \textbf{Hierarchy and Scale:} The TOR/Tier-2/Tier-1 architecture supports tens of thousands of hosts as a unified system.
    \item \cite{13778} \textbf{Capacity vs.\ Redundancy:} Interconnecting switches via Clos/Fat-Tree fabrics provides path diversity and supports multi-path routing.
    \item \cite{13780} \textbf{Congestion and Latency:} Low-latency DCNs require ECN-driven congestion control rather than TCP’s loss-based approach.
\end{itemize}

\subsubsection{Key Insights / Takeaways}

\begin{itemize}
    \item \textbf{Complexity at the Edge:} Although the underlying DCN architecture is complex, it is engineered so that hosts experience a simple, uniform, high-speed network.
    \item \cite{13779} \textbf{Software-Defined Infrastructure:} Centralized SDN control is critical for scale, automation, and global optimization.
    \item \cite{13777,13779} \textbf{Networking as a Cost Factor:} Efficient multi-path routing and congestion control directly influence operational cost and performance.
\end{itemize}


\section{Retrospective: A Day in the Life of a Web Page Request (Section 6.7)}
\label{sec:day-in-life}

In this section, we take an integrated, holistic journey down and up the protocol stack by following a single Web page request from a user's browser. As emphasized in the textbook, even a ``simple’’ request to download a Web page involves \emph{many (many!)} protocols working together across all layers of the Internet architecture. This retrospective synthesizes DHCP, DNS, ARP, IP, TCP, and HTTP into one continuous narrative of encapsulation, addressing, and hop-by-hop forwarding. Figure~6.32 in the textbook depicts the setting: a student, Bob, plugs his laptop into an Ethernet switch at his school and downloads a page such as \url{www.google.com}. As we will see, there is a great deal happening ``under the hood.'' (A Wireshark lab at the end of the chapter analyzes similar traces.)

\subsection{Concept Overview: The Synthesis of the Protocol Stack}

Typing a URL into a browser triggers an extended sequence of protocol interactions. The goals of this retrospective are to (i) trace the full lifecycle of a Web request from application layer down to the physical link and back, and (ii) clearly identify the roles played by DHCP, DNS, ARP, IP, TCP, and HTTP.

We assume a client host~C wants to fetch a Web page from a remote server~S at \texttt{www.somecollege.edu}. The host is connected to a local Ethernet switch, which is connected to a first-hop router~R1. As in the textbook scenario, the router is connected to an ISP; for example, Comcast may provide the DNS service for the school, so the DNS server resides \emph{not} on the school's LAN but within the ISP network.

\subsection{Phase I: Initial Configuration and Name Resolution (L7, L4, L3, L2)}

Before C can send any application-layer request, it must acquire network configuration and learn the server's IP address. This involves DHCP and DNS.

\subsubsection{Initial Network Configuration via DHCP}

If C is newly connected to the LAN, it must obtain an IP address and other parameters through the Dynamic Host Configuration Protocol (DHCP), an application-layer protocol running over UDP.

\paragraph{Four-Step DHCP Process}
\begin{enumerate}
    \item \textbf{DHCP Discover (Broadcast):}
          C broadcasts a DHCP Discover message to \texttt{255.255.255.255} with destination MAC \texttt{FF:FF:FF:FF:FF:FF}. Because C has no IP address, the Discover is encapsulated in a UDP segment within an IP datagram whose source IP is \texttt{0.0.0.0}. This datagram is then carried inside an Ethernet frame and sent into the LAN.

    \item \textbf{DHCP Offer (Unicast or Broadcast):}
          A DHCP server receives the broadcast and replies with a DHCP Offer proposing an IP address $IP_C$ and associated configuration.

    \item \textbf{DHCP Request (Broadcast):}
          If multiple offers arrive, C picks one and broadcasts a DHCP Request message, specifying the chosen server's offer.

    \item \textbf{DHCP ACK (Unicast):}
          The server responds with a DHCP ACK that includes:
          \begin{itemize}
              \item $IP_C$ (the leased IP address),
              \item subnet mask,
              \item default gateway (router R1),
              \item the IP address of the DNS server (here provided by Comcast in the textbook scenario).
          \end{itemize}
\end{enumerate}

\paragraph{Intuition}
This process resembles a newcomer entering an office environment and asking an administrator (DHCP) for their desk number, campus mailbox rules (subnet mask), manager (default gateway), and directory services (DNS).

\subsubsection{DNS Query: Resolving the Web Server's IP}

Before sending an HTTP request, the browser must obtain the server's IP $IP_S$ from the Domain Name System (DNS).

\paragraph{Client-Side Operation}
\begin{itemize}
    \item The browser extracts the hostname from the URL and hands it to the DNS client.
    \item The DNS client constructs a DNS query message (application layer).
    \item The query is placed into a UDP segment (transport layer) with destination port~53.
    \item The UDP segment is encapsulated into an IP datagram (network layer) with destination equal to the DNS server's IP (provided by Comcast in the text's example).
\end{itemize}

At this point, the datagram must be sent to R1 for forwarding toward the DNS server.

\subsection{Phase II: Establishing Link-Layer Connectivity (L2)}

\subsubsection{ARP: Resolving the Router’s MAC Address}

To send the DNS query to R1, C must encapsulate the IP datagram in a link-layer frame addressed to R1's MAC address.

\paragraph{Technical Mechanism}
\begin{enumerate}
    \item C checks its ARP cache; if $MAC_{R1}$ is absent, it broadcasts an ARP query frame.
    \item The ARP frame uses destination MAC \texttt{FF:FF:FF:FF:FF:FF} and asks: ``Who has $IP_{R1}$?''
    \item Every adapter on the LAN receives the broadcast; R1 recognizes its IP and replies with an ARP Response containing $MAC_{R1}$.
    \item C updates its ARP cache and can now send frames addressed to $MAC_{R1}$.
\end{enumerate}

\paragraph{Key Principle}
ARP highlights the separation between network-layer addressing (IP) and link-layer addressing (MAC). C \emph{never} needs the MAC address of the remote Web server; only the MAC of the next-hop router at each link.

\subsection{Phase III: Request and Response (L7, L4, L3, L2)}

With DNS and ARP completed, C can resolve the server address and initiate TCP/HTTP.

\subsubsection{DNS Resolution and the TCP Handshake}

\paragraph{DNS Response}
The DNS server returns $IP_S$ to C, completing name resolution.

\paragraph{TCP Connection Setup}
C creates a TCP SYN segment destined for S:
\begin{itemize}
    \item The SYN is encapsulated in an IP datagram with destination $IP_S$.
    \item The IP datagram is encapsulated in a link-layer frame addressed to $MAC_{R1}$.
    \item R1 forwards the datagram hop-by-hop across the Internet.
    \item The server replies with a SYNACK; C responds with an ACK, completing the three-way handshake.
\end{itemize}

\subsubsection{HTTP Transaction and Data Transfer}

\paragraph{Technical Mechanism}
\begin{enumerate}
    \item C sends an HTTP \texttt{GET} request for \texttt{/index.html} over the established TCP connection.
    \item S processes the GET and retrieves the HTML file.
    \item S sends the file as a stream of TCP segments; TCP handles reliable, in-order delivery and congestion control.
    \item Segments are carried in IP datagrams; datagrams are encapsulated in link-layer frames and forwarded across each hop.
\end{enumerate}

\paragraph{Router Encapsulation/De-encapsulation Example}

Consider an intermediate router~R2:
\begin{itemize}
    \item \textbf{Inbound at R2:}
          R2 receives a link-layer frame, removes the frame header, checks for errors, removes the IP header, inspects the destination address, and consults its forwarding table.
    \item \textbf{Outbound at R2:}
          R2 decrements the IP TTL, determines the next hop (R3), uses ARP if needed to obtain $MAC_{R3}$, constructs a new frame with a fresh link-layer header, and transmits it.
\end{itemize}

\[
    \text{Frame}_{in} \to \text{Datagram} \to \text{Datagram (TTL-1)} \to \text{Frame}_{out}
\]
\[
    MAC_{in} \to IP \to IP \to MAC_{out}
\]

This process is repeated at every router: full link-layer decapsulation and re-encapsulation, while the IP datagram remains mostly unchanged end-to-end.

\subsubsection{Final Delivery}

The TCP segments arrive at C:
\begin{itemize}
    \item TCP reassembles the byte stream.
    \item The application layer (browser) processes the HTML.
    \item The Web page appears to the user.
\end{itemize}

\subsection{Summary and Reinforcement}

\subsubsection{Key Terms and Definitions}

\begin{itemize}
    \item \textbf{Link-Local Discovery:} Processes such as DHCP and ARP that allow a host to obtain configuration and determine the MAC address of the next hop.
    \item \textbf{Encapsulation/De-encapsulation:} Headers added as data moves down the stack and removed as data moves up; routers decapsulate only to the network layer.
    \item \textbf{Three-Way Handshake:} TCP's SYN--SYNACK--ACK procedure that establishes a reliable connection.
    \item \textbf{IP Datagram vs.\ Link-Layer Frame:} The IP datagram's addresses remain constant end-to-end; link-layer frames change at every hop.
\end{itemize}

\subsubsection{Core Relationships}

\begin{itemize}
    \item \textbf{Interdependence of Layers:} Application-layer protocols depend on transport-layer services, which depend on IP routing, which requires link-layer addressing.
    \item \textbf{Broadcast vs.\ Unicast:} DHCP Discover and ARP queries use broadcast; DNS, TCP, and HTTP use unicast once addressing is resolved.
\end{itemize}

\subsubsection{Key Insights / Takeaways}

\begin{itemize}
    \item \textbf{TTL is the Only Changing IP Header Field Hop-by-Hop:} Each router decrements the TTL; otherwise the IP datagram stays intact.
    \item \textbf{Layering Abstraction Works:} HTTP is oblivious to whether frames travel over Ethernet, Wi-Fi, MPLS, or a mixture.
    \item \textbf{Initial Setup Dominates Complexity:} DHCP, ARP, and DNS often require more protocol exchanges than the ensuing HTTP data transfer.
\end{itemize}



% template
\begin{itemize}
    \item X
          \begin{enumerate}
              \item Y
          \end{enumerate}
\end{itemize}

\end{document}